% Appendix: baseline adaptations and settings.
\section{Comparison Systems}
\label{app:baselines}

All comparison systems (B1--B4, X1, X2) use the same cast list, character-page template,
generator, and intended character-page budget:
the shared sections and sentence limits (appearance 2--4, personality 2--6, history 3--12), at most five
quotes, and a shared citation format (Table~\ref{tab:baselines}). For systems supplied with
passages, citations to paragraphs outside their input are removed, and characters absent
from the eligible text receive no page but remain in the evaluation denominators. Settings were chosen on the development novel
(\emph{Middlemarch}) and fixed before the Anne runs. \sys itself is not given the cast: it
discovers entities and also publishes pages for other characters, relationships,
codex entries, and a timeline. The shared template therefore controls the evaluated cast's
character pages, while the complete artifacts differ in scope (Table~\ref{tab:kinds}).
The exhibited comparison prompts request 2--4 appearance sentences; the archived \sys
appearance prompt requests 3--4 (Appendix~\ref{app:prompts}). The template budget therefore
does not imply identical prompt wording or verified enforcement of output length.


\begin{table}[t]
\centering\small
\begin{tabular}{@{}lp{6.5cm}@{}}
\toprule
 & \nwhead{Adaptation and settings} \\
\midrule
B1 & Every paragraph of volumes $1,\ldots,t$, labeled with its identifier, precedes the task so that pages at a cutoff share a cacheable prefix. Prompts exceeding the model's context limit would be marked unsupported rather than truncated; none exceeded it. \\
B2 & Prefix-only index of consecutive paragraphs within a chapter (150 words). Okapi BM25 ($k_1{=}1.5$, $b{=}0.75$) and BGE-M3 cosine rankings for the character name alone and with seven facet words (appearance, personality, family, relations, goals, past, speech), fused by reciprocal rank ($k{=}60$); chunks added in fused order up to 20{,}000 tokens of retrieved text, then restored to reading order. Budgets count retrieved text only, estimated as characters divided by four, so a call's provider-counted input, which adds the prompt, can exceed them. Chunk size (150 or 300 words) and budget (6{,}000, 12{,}000 or 20{,}000 tokens) were chosen at $t{=}2$ of the development novel. \\
B3 & LightRAG 1.5.7, default chunking (1{,}200 tokens) and its default extraction prompt and limits (at most 100 rows and 40 entities per chunk), mode \texttt{mix}, context capped at the same 20{,}000 tokens, reranking off. The index at cutoff $t$ is the index at $t{-}1$ plus volume $t$, so no index sees a later volume. LightRAG supplies retrieval only: its retrieved entities, relations and chunks are passed to the shared generator, and its own extraction calls use the same model. \\
B4 & The \sys graph built from all five volumes. For each character, the facts visible at $t$ and their evidence paragraphs of volumes $\le t$, up to the same budget, are passed to the shared generator in place of \sys's structured rendering (graph \texttt{anne@v2}). \\
X1 & The shared prompt with the work's title, the character's name and no passages, with an instruction to write what is known of the character up to volume $t$; citations are empty by construction. \\
X2 & Every paragraph of volumes 1--5 regardless of $t$, with the instruction to write for a reader of volume $t$ and reveal nothing later. \\
\bottomrule
\end{tabular}
\caption{Adaptations of the comparison systems.}
\label{tab:baselines}
\end{table}
